\subsection{9.45: a coordinate test for orthogonal lattice bases}
\label{cand:9.45}
For \(a\in\Q^n\), let \(S(a)=\Z^n+\Z a\).
Subtracting an integer vector, write
\[
 a=r/m,\quad 0\le r_j<m,\quad \gcd(m,r_1,\ldots,r_n)=1.
\]
The following finite test uses \(O(mn)\) integer arithmetic operations;
no polynomial bound in the binary length of \(m\) is asserted.

For each \(1\le k<m\), center every coordinate of \(kr\bmod m\) in
\([-m/2,m/2]\). At a half-coordinate retain both signs; discard a coset
with more than two such ties. Among the resulting at most four vectors
\(w\), retain those with first nonzero coordinate positive which satisfy,
with \(Q(w)=\sum_jw_j^2\),
\[
 Q(w)\mid mw_j\quad\hbox{for every }j,\qquad
 Q(w)\mid\sum_jr_jw_j.                              \tag{**}
\]
Call their list \(\mathcal W\), and let \(z\) count the zero coordinates
of \(r\).

\begin{proposition}
The lattice \(S(a)\) has an orthogonal basis if and only if
\(z+|\mathcal W|=n\). In that case a basis is
\(\{e_j:r_j=0\}\cup\{w/m:w\in\mathcal W\}\).
\end{proposition}
\begin{proof}
For a Euclidean lattice \(L\), a vector \(b\ne0\) spans an orthogonal
rank-one direct summand exactly when \(c_b=b/\|b\|^2\in L^*\):
subtract the integer multiple \((x\cdot c_b)b\) from each \(x\in L\).
Call such a vector splitting. For splitting vectors \(b,v\), both
\[
 \frac{b\cdot v}{\|b\|^2},\qquad
 \frac{b\cdot v}{\|v\|^2}
\]
are integers. Their product is the squared cosine. If the lines differ,
it is less than one and hence zero; if they coincide, the ratio and its
reciprocal are integers, so \(v=\pm b\). All splitting lines are
therefore mutually orthogonal. There are at most \(n\), and if there are
\(n\), the orthogonal expansion of every \(x\in L\) has integer
coefficients, proving that their generators form a basis.

Now \(L=S(r/m)\) contains \(\Z^n\), so \(L^*\subseteq\Z^n\).
For a splitting vector write \(c=c_b\), \(d=\|c\|^2\).
Then \(c\in\Z^n\), \(d\) is a positive integer, and \(b=c/d\).
If \(d=1\), \(b=\pm e_j\), which splits precisely when \(r_j=0\).
If \(d\ge2\), then \(2|c_j|\le d\) in every coordinate, so
\(|b_j|\le1/2\). If \(h\) coordinates attain equality, then
\(hd^2/4\le d\), giving \(h\le2\).
Moreover \(\|b\|^2=1/d<1\), so \(b\notin\Z^n\).
The gcd condition shows that \(r/m\) has order exactly \(m\) modulo
\(\Z^n\). Thus every remaining splitting vector, up to its unique
first-positive sign, occurs in the centered enumeration.

For a listed \(w\), \(b=w/m\) lies in \(L\), and \(c_b=mw/Q(w)\).
Membership in \(L^*\) says that its coordinates and its scalar product
with \(r/m\) are integers, which is exactly (**).
Different cosets or tie choices give different vectors. The test
therefore lists all splitting lines exactly once, proving the result.
\end{proof}

Half-coordinate signs matter: for \(a=(1/6,1/6)\), the vector
\((3,-3)\) passes (**) while \((3,3)\) fails. For
\(a=(1/2,1/2,1/2)\), there are no splitting lines.
A three-dimensional positive example is \(a=(1,5,2)/6\), with basis
\[
 (1,-1,2)/6,\qquad (1,-1,-1)/3,\qquad (1,1,0)/2.
\]
The archive also proves that when \(m=p^k\), at most two coordinates
of \(r\) can be nonzero in an orthogonal example. Indeed an integral
orthogonal basis matrix \(C\) of \(L^*\) has determinant \(p^k\),
Smith form \(\operatorname{diag}(1,\ldots,1,p^k)\), and diagonal
\(CC^t\) whose entries are powers of \(p\).
The rank inequality modulo \(p\) gives
\(\operatorname{rank}(CC^t)\ge n-2\).
All but at most two row norms are therefore one, forcing the other
residue coordinates to vanish.

Bardakov already solved the two-dimensional case and gave a
higher-dimensional sufficient condition \cite{bardakov-lattices}.
Normalized-dual arguments and orthogonal decomposition are prior
methods; compare \cite{cgg-orthogonality,hemkemeier-vallentin}.
The candidate contribution is this particular cyclic-coset criterion
in arbitrary dimension. The archive
\artifact{research/9.45-proof.md} contains the prime-power classification;
\artifact{results/9.45-summary.json} binds independent finite controls.
